\chapter{Conclusiones y Recomendaciones}

\section{Conclusiones}
Se desarrolló y evaluó un prototipo académico de clasificación de la desnutrición crónica infantil con microdatos de la ENSANUT 2018. Las conclusiones siguientes corresponden a cada objetivo específico y se sustentan en los resultados del Capítulo~III.

\noindent\hspace*{\parindent}\textbf{Objetivo 1: análisis y preparación de datos.} La base original tenía \num{20510} registros y 350 variables; \num{2017} registros (\SI{9.8}{\percent}) carecían de la variable objetivo y se excluyeron en lugar de imputarse, con lo que se conservaron \num{18493} registros y 114 predictores. Los predictores no incluyen talla ni peso actuales del niño, pero sí dos variables de diseño muestral (\texttt{factor\_expansion} y \texttt{estrato}) que pueden condicionar los resultados. Además, 84 de los 114 predictores tenían valores ausentes (32 con más del \SI{10}{\percent}), que se imputaron antes de la partición.

\noindent\hspace*{\parindent}\textbf{Objetivo 2: diseño de la arquitectura y del protocolo.} La solución se organizó en cuatro componentes (pipeline, artefactos, API y tablero). El protocolo fijó una partición 80/20 estratificada con semilla 42, un preprocesamiento dentro del \textit{Pipeline}, la selección por PR-AUC y un umbral de decisión de 0.42. Se definieron 8 requerimientos funcionales y 7 no funcionales con criterios de aceptación.

\noindent\hspace*{\parindent}\textbf{Objetivo 3: entrenamiento, comparación y evaluación.} Random Forest obtuvo el mayor PR-AUC entre los cuatro modelos (0.411). Con el modelo final, en el conjunto de prueba nacional (n = \num{3699}), la sensibilidad fue de 0.693, la precisión de 0.335, el F1 de 0.452, la exactitud balanceada de 0.621, el ROC-AUC de 0.677 y el PR-AUC de 0.409. En Tungurahua (n = 112) la sensibilidad fue de 0.686 con intervalo de \SI{95}{\percent} de 0.533 a 0.833. El desempeño es moderado y comparable con el reportado con datos de la ENDI (AUC de 0.62 a 0.66), y el modelo presenta sobreajuste (exactitud balanceada de 0.819 en entrenamiento frente a 0.621 en prueba).

\noindent\hspace*{\parindent}\textbf{Objetivo 4: API y tablero.} Se implementó una API REST con nueve servicios, que valida las 114 variables de entrada, y un tablero web que consume la API para consultar indicadores, métricas y predicciones.

\noindent\hspace*{\parindent}\textbf{Objetivo 5: validación funcional y limitaciones.} Las 8 pruebas automáticas y 15 de los 17 casos de prueba se cumplieron (fallaron el envío de un Excel a la API y la variación de la puntuación del formulario individual), y la latencia mediana fue de 29 ms en ejecución local. Estas pruebas verifican el funcionamiento del software y no la validez predictiva ni clínica. Se identificó que el nivel de riesgo contradecía la clase asignada (corregido: BAJO coincide ahora con la clase 0), que el resumen del tablero usa predicciones sobre datos de entrenamiento, que la puntuación del formulario individual casi no varía con los datos porque el resto de las variables se imputa, y que la API no tiene autenticación.

\subsection{Limitaciones}
\begin{itemize}
  \item \textbf{Representatividad territorial:} Tungurahua aporta \num{531} registros (112 en la prueba), por lo que sus estimaciones son imprecisas y no permiten resultados por cantón o parroquia. Las prevalencias reportadas no están ponderadas.
  \item \textbf{Diseño transversal:} el modelo estima la condición observada; no anticipa casos futuros.
  \item \textbf{Desbalance:} la clase de interés es el \SI{24.7}{\percent}; la precisión es baja (una de cada tres alertas es un caso real).
  \item \textbf{Posible fuga o sesgo por variables de diseño:} \texttt{factor\_expansion} y \texttt{estrato} figuran entre las variables más importantes.
  \item \textbf{Imputación previa a la partición:} los faltantes se imputaron con mediana o moda sobre todos los registros antes de dividir los datos, y el peso y la talla al nacer tienen más de la mitad de sus valores imputados. El efecto sobre las métricas no se cuantificó.
  \item \textbf{Selección del modelo:} la comparación de los cuatro modelos base se hizo sobre la partición de prueba, por lo que el PR-AUC de prueba puede ser levemente optimista.
  \item \textbf{Validación externa:} no existe. La ENDI se exploró pero no se usó para validar.
  \item \textbf{Seguridad y reproducibilidad:} la API no autentica usuarios y el CORS está abierto. Los resultados dependen de las versiones de librerías indicadas.
  \item \textbf{Alcance clínico:} el prototipo es una herramienta académica de apoyo analítico; no es un sistema de diagnóstico ni una solución validada para decisiones médicas.
  \item \textbf{Calibración y subgrupos:} las probabilidades no están calibradas y la sensibilidad varía según el área y la edad (Capítulo~III), por lo que un umbral único no es adecuado para un uso aplicado.
\end{itemize}

\subsection{Aportes del trabajo}
\begin{itemize}
  \item \textbf{Procedimiento documentado y corregido:} se detectó y corrigió una fuga en la versión inicial (\num{2017} etiquetas imputadas) y se dejó el procedimiento reproducible, con semilla fija y código disponible (Anexo~I).
  \item \textbf{Evaluación transparente:} se reportaron métricas por clase, intervalos de confianza, sobreajuste, calibración y desempeño por subgrupos, en lugar de una exactitud global.
  \item \textbf{Evaluación para Tungurahua:} se cuantificó el desempeño en la provincia con sus límites de representatividad.
  \item \textbf{Prototipo con pruebas:} se entregó una API con validación de entradas, casos de prueba y una lista explícita de defectos encontrados.
\end{itemize}

\subsection{Lectura conjunta de los análisis complementarios}
Los análisis adicionales del Capítulo~III matizan el desempeño global. Primero, las probabilidades del modelo no están calibradas (error de Brier de 0.205 frente a 0.186 de la prevalencia), por lo que sirven para ordenar casos y no como riesgo absoluto. Segundo, con un umbral único la sensibilidad cambia mucho según el subgrupo: 0.880 en el área rural frente a 0.525 en la urbana, y 0.892 entre los 12 y 23 meses frente a 0.505 entre los 36 y 59 meses. Tercero, excluir las variables de diseño muestral no modifica el ROC-AUC (0.679 frente a 0.677), de modo que el desempeño no depende de ellas. Cuarto, la elección de Random Forest sobre la regresión logística es una preferencia débil: el ROC-AUC de la regresión logística (0.667) y su PR-AUC (0.391) son cercanos a los del bosque (0.680 y 0.411).



\section{Recomendaciones y Trabajos Futuros}
\subsection{Recomendaciones inmediatas}
\begin{itemize}
  \item \textbf{Corregir el resumen del tablero:} calcular los indicadores con los datos observados o con el conjunto de prueba, no con predicciones sobre datos de entrenamiento.
  \item \textbf{Corregir la lectura de archivos en la API:} aceptar Excel o rechazarlo con un error controlado.
  \item \textbf{Decidir el futuro del formulario individual:} mantenerlo solo como ilustración o reemplazarlo por una carga de un registro completo.
  \item \textbf{Proteger la API:} agregar autenticación y restringir el origen de las solicitudes antes de cualquier despliegue abierto.
\end{itemize}

\subsection{Trabajos futuros}
\begin{itemize}
  \item Rehacer la imputación dentro del \textit{Pipeline} a partir de la base original, para cuantificar su efecto, y evaluar con ponderación muestral.
  \item Validar con una fuente independiente, como la ENDI, con una variable objetivo comparable.
  \item Calibrar las probabilidades y analizar la interpretabilidad con métodos como SHAP.
  \item Contrastar el desempeño con profesionales de salud antes de cualquier uso aplicado.
\end{itemize}
